% !PRINTED 146
\weberpage{146}
%<<CONT
\weq{8}{\begin{gathered}
 C_0\alpha_1^{n-1}+C_1\alpha_1^{n-2}+\cdots+C_{n-2}\alpha_1+(C_{n-1}-S)=0\\
 C_0\alpha_2^{n-1}+C_1\alpha_2^{n-2}+\cdots+C_{n-2}\alpha_2+(C_{n-1}-S)=0\\
 \makebox[30em]{\dotfill}\\
 C_0\alpha_n^{n-1}+C_1\alpha_n^{n-2}+\cdots+C_{n-2}\alpha_n+(C_{n-1}-S)=0.
\end{gathered}}
Betrachten wir dies System als ein System homogener linearer Gleichungen mit den $n$ Unbekannten
\[
 C_0,C_1\ldots C_{n-2},C_{n-1}-S,
\]
so ist seine Determinante
\[
\begin{vmatrix}
 \alpha_1^{n-1},&\alpha_1^{n-2}&\ldots&\alpha_1,&1\\
 \alpha_2^{n-1},&\alpha_2^{n-2}&\ldots&\alpha_2,&1\\
 \multicolumn{5}{c}{\makebox[12em]{\dotfill}}\\
 \alpha_n^{n-1},&\alpha_n^{n-2}&\ldots&\alpha_n,&1
\end{vmatrix},
\]
die nach \S.\,22, Formel (12) gleich dem Product aller Differenzen $\alpha_1-\alpha_2,\alpha_1-\alpha_3\ldots$ ist, von Null verschieden, und folglich ist nach dem Satz II in \S.\,24:
\[
 C_0=0,\quad C_1=0\ldots\quad C_{n-2}=0
\]
und
\weq{9}{S=C_{n-1},}
worin der zu beweisende Fundamentalsatz enthalten ist.

Zu demselben Resultat gelangt man auch durch den Satz II, \S.\,30; denn da die Function $(n-1)^{\text{ten}}$ Grades von $x$
\[
 C_0x^{n-1}+C_1x^{n-2}+\cdots+C_{n-2}x+C_{n-1}-S
\]
für $x=\alpha_1,\alpha_2\ldots\alpha_n$, also für mehr als $n-1$ Werthe verschwindet, so müssen ihre Coëfficienten alle Null sein.

Der Beweis, den wir hier für das Fundamentaltheorem im Anschluss an \gesperrt{Cauchy}\footnote[1]{\gesperrt{Cauchy}, Exercices de mathématiques, $4^{\text{ème}}$ année.} gegeben haben, bietet zugleich ein Mittel, in besonderen Fällen den Ausdruck einer symmetrischen Function durch die Grundfunctionen wirklich zu berechnen. Dieselbe Möglichkeit bieten auch die anderen Beweise, die für das Theorem bekannt sind. Wir wollen noch einen zweiten Beweis hier mittheilen, der zu einer oft einfacheren Berechnungsart führt\footnote[2]{\gesperrt{Waring}, Meditationes algebraicae. \gesperrt{Gauss}, Demonstratio nova altera theorematis omnem functionem algebraicam rationalem integram unius variabilis in factores reales primi vel secundi gradus resolvi posse. Werke, Bd.\ III, S.\ 36.}.
